\chapter{Cơ sở lý thuyết và công nghệ nền tảng}

\section{Hệ thống quản lý tài liệu bảo mật}

Hệ thống quản lý tài liệu (Document Management System - DMS) là hệ thống hỗ trợ lưu trữ, tổ chức, tìm kiếm, chia sẻ và kiểm soát vòng đời tài liệu. Trong môi trường doanh nghiệp, một DMS quản lý cả file lẫn metadata, phiên bản, người sở hữu, trạng thái phê duyệt và lịch sử thao tác.

Với tài liệu nhạy cảm, DMS cần thêm các yêu cầu bảo mật. Một tài liệu không nên chỉ được bảo vệ bằng vị trí lưu trữ, mà phải được bảo vệ bằng chính sách truy cập nhất quán ở backend, đặc biệt với các thao tác như preview, stream, presign download, export evidence hoặc truy vấn audit. Nếu hệ thống chỉ ẩn nút thao tác ở frontend trong khi backend vẫn cho gọi API, chính sách bảo mật không còn đáng tin cậy.

DocVault được thiết kế theo quan điểm backend là nơi thực thi chính sách cuối cùng. Frontend có thể hiển thị hoặc vô hiệu hóa thao tác để cải thiện trải nghiệm người dùng, nhưng quyết định cho phép hoặc từ chối vẫn thuộc về metadata-service, document-service, gateway và audit-service.

\section{Khảo sát các hệ thống tương tự}

Trên thị trường đã có nhiều hệ thống quản lý tài liệu, từ giải pháp thương mại đến mã nguồn mở. Việc đối chiếu với các hệ thống này giúp xác định khoảng trống mà DocVault hướng tới. Bảng~\ref{tab:related-work} so sánh DocVault với ba đại diện tiêu biểu trên các tiêu chí liên quan trực tiếp đến phạm vi đề tài.

\begin{itemize}
  \item \textbf{Microsoft SharePoint}: nền tảng quản lý tài liệu và cộng tác thương mại, mạnh về tích hợp hệ sinh thái Microsoft 365, phân quyền và compliance; tuy nhiên là dịch vụ đóng, chủ yếu vận hành trên đám mây của nhà cung cấp và khó tùy biến sâu chính sách truy cập ở mức mã nguồn.
  \item \textbf{Alfresco}: hệ thống ECM mã nguồn mở (có bản thương mại), hỗ trợ vòng đời tài liệu, phân quyền và audit; kiến trúc nghiêng về monolith/Java truyền thống và không đặt trọng tâm vào quy trình DevSecOps đi kèm.
  \item \textbf{Nextcloud}: giải pháp self-hosted phổ biến cho lưu trữ và chia sẻ file, dễ triển khai; mạnh về chia sẻ và đồng bộ nhưng kiểm soát theo classification, DLP và audit tamper-evident không phải trọng tâm mặc định.
\end{itemize}

\begin{table}[H]
\centering
\small
\caption{So sánh DocVault với một số hệ thống quản lý tài liệu tiêu biểu}
\label{tab:related-work}
\begin{tabularx}{\textwidth}{p{0.26\textwidth}YYYY}
\toprule
\textbf{Tiêu chí} & \textbf{SharePoint} & \textbf{Alfresco} & \textbf{Nextcloud} & \textbf{DocVault} \\
\midrule
Mô hình triển khai & Cloud đóng & Self-hosted & Self-hosted & Self-hosted \\
RBAC + ACL deny-overrides & Có & Có & Một phần & Có \\
Phân loại + kiểm soát theo classification & Có & Một phần & Hạn chế & Có \\
DLP nâng classification tự động & Một phần & Hạn chế & Hạn chế & Có \\
Audit tamper-evident (hash-chain) & Hạn chế & Hạn chế & Hạn chế & Có \\
Compliance boundary (metadata-only) & Hạn chế & Hạn chế & Hạn chế & Có \\
DevSecOps pipeline đi kèm & Không & Không & Không & Có \\
\bottomrule
\end{tabularx}
\end{table}

Điểm khác biệt của DocVault không nằm ở việc thay thế các hệ thống trên về quy mô tính năng, mà ở cách gộp ba yếu tố trong cùng một đề tài: (1) kiểm soát truy cập nhiều lớp theo classification với ranh giới compliance metadata-only, (2) audit tamper-evident bằng hash-chain, và (3) đưa toàn bộ quy trình phát hành vào pipeline DevSecOps có bằng chứng. Hướng tiếp cận này bám sát mục tiêu học thuật của đề tài: minh họa cách bảo mật được đưa xuyên suốt từ thiết kế ứng dụng đến quy trình triển khai, thay vì chỉ dựng một DMS đầy đủ tính năng.

\section{Kiến trúc microservices}

Microservices là cách tổ chức hệ thống thành nhiều service nhỏ, mỗi service sở hữu một năng lực nghiệp vụ rõ ràng và có thể phát triển, kiểm thử, triển khai độc lập ở mức tương đối. Trong DocVault, kiến trúc này được dùng để tách các bounded context:

\begin{itemize}
  \item Gateway chịu trách nhiệm xác thực, routing và truyền ngữ cảnh.
  \item Metadata-service sở hữu metadata, ACL, trạng thái và policy.
  \item Document-service sở hữu upload/download file và tích hợp object storage.
  \item Workflow-service sở hữu state machine phê duyệt.
  \item Audit-service sở hữu audit trail và hash-chain.
  \item Notification-service nhận thông báo nghiệp vụ trong môi trường demo/dev.
\end{itemize}

Cách tách này giúp giảm sự phụ thuộc giữa lưu trữ file và chính sách truy cập. File thật nằm ở MinIO/S3, còn metadata và quyền truy cập nằm ở service khác. Khi người dùng muốn tải file, document-service không tự ý quyết định quyền mà dựa vào policy/grant token do metadata-service cấp hoặc xác minh lại.

\section{Xác thực và phân quyền}

\subsection{Keycloak, OIDC và JWT}

Keycloak là hệ thống quản lý định danh và truy cập mã nguồn mở, hỗ trợ các chuẩn phổ biến như OAuth 2.0 và OpenID Connect \cite{keycloak_docs}. Trong DocVault, Keycloak được sử dụng để quản lý người dùng, vai trò, nhóm và phát hành token. Sau khi đăng nhập, người dùng nhận JWT chứa thông tin định danh và các claim cần thiết như user id, username, roles và groups.

Gateway và các service phía sau dùng JWT để xác định actor của request. Gateway là entry point chính, xác thực token và truyền các header ngữ cảnh như \code{x-user-id}, \code{x-roles}, \code{x-groups} xuống các service nội bộ. Nhờ đó service nghiệp vụ không phụ thuộc trực tiếp vào phiên đăng nhập frontend mà vẫn đủ thông tin để thực thi chính sách.

\subsection{RBAC}

RBAC là mô hình phân quyền theo vai trò. Thay vì gán quyền trực tiếp cho từng người dùng, hệ thống gán quyền cho vai trò, sau đó người dùng được gán một hoặc nhiều vai trò. Trong DocVault, các vai trò chính gồm:

\begin{itemize}
  \item \code{viewer}: đọc tài liệu đã xuất bản nếu policy cho phép.
  \item \code{editor}: tạo tài liệu, chỉnh sửa metadata, upload file và gửi duyệt.
  \item \code{approver}: xem và phê duyệt hoặc từ chối tài liệu đang chờ duyệt.
  \item \code{compliance\_officer}: xem audit, metadata và evidence trong phạm vi compliance, nhưng không được truy cập nội dung file.
  \item \code{admin}: quản trị hệ thống và có quyền rộng nhất.
\end{itemize}

RBAC phù hợp để kiểm soát quyền ở mức route hoặc chức năng, nhưng chưa đủ để mô hình hóa mọi trường hợp truy cập theo từng tài liệu. Vì vậy, DocVault bổ sung ACL.

\subsection{ACL và deny-overrides}

ACL cho phép gán rule truy cập ở mức tài liệu. Một rule có thể áp dụng cho user, role, group hoặc tất cả người dùng. Mỗi rule có permission như \code{READ}, \code{DOWNLOAD}, \code{WRITE}, \code{APPROVE} và effect là \code{ALLOW} hoặc \code{DENY}.

DocVault sử dụng chiến lược deny-overrides: nếu có rule \code{DENY} phù hợp, kết quả cuối cùng là từ chối, ngay cả khi có rule \code{ALLOW} khác. Đây là chiến lược an toàn hơn cho dữ liệu nhạy cảm, vì một lệnh cấm rõ ràng không bị ghi đè bởi quyền rộng hơn.

\section{Phân loại tài liệu và kiểm soát nội dung}

DocVault sử dụng bốn mức phân loại tài liệu:

\begin{itemize}
  \item \code{PUBLIC}: tài liệu công khai trong phạm vi người dùng đã xác thực.
  \item \code{INTERNAL}: tài liệu nội bộ.
  \item \code{CONFIDENTIAL}: tài liệu mật, cần kiểm soát preview/download nghiêm ngặt hơn.
  \item \code{SECRET}: tài liệu rất nhạy cảm, chỉ một số vai trò/đối tượng được phép truy cập.
\end{itemize}

Mức phân loại không dừng ở metadata hiển thị mà còn ảnh hưởng đến chính sách truy cập. Chẳng hạn, tài liệu \code{CONFIDENTIAL} hoặc \code{SECRET} có thể không trả về presigned URL trực tiếp mà phải đi qua endpoint stream, để document-service áp thêm kiểm soát như watermark posture và tái xác thực grant.

\section{Audit, hash-chain và compliance evidence}

Audit log ghi lại các sự kiện quan trọng trong hệ thống: tạo tài liệu, upload file, submit, approve, reject, download authorized/denied, DLP detected, malware blocked, retention archive và các thao tác compliance. Một audit log hữu ích thì ngoài việc có dữ liệu còn phải cho phép kiểm tra tính toàn vẹn.

DocVault sử dụng hash-chain cho audit events. Mỗi event có hash phụ thuộc vào nội dung event hiện tại và hash của event trước đó. Nếu một event cũ bị sửa trực tiếp trong storage, chuỗi hash phía sau sẽ không còn khớp và verify-chain có thể phát hiện bất thường. Cơ chế này không biến audit log thành lưu trữ bất biến tuyệt đối, nhưng tạo tính tamper-evident phù hợp với phạm vi đồ án.

Bên cạnh audit query, hệ thống còn cung cấp evidence packet. Evidence packet gom metadata, version checksum, workflow history, retention evidence, trạng thái verify-chain và các audit event liên quan đến một tài liệu. Đáng chú ý là evidence packet không chứa nội dung file, presigned URL, grant token hoặc object key nhạy cảm.

\section{DevSecOps}

DevSecOps là cách tiếp cận tích hợp bảo mật vào quy trình phát triển và vận hành. Thay vì kiểm tra bảo mật ở cuối dự án, DevSecOps đưa các kiểm soát vào nhiều điểm trong pipeline: kiểm tra mã nguồn, kiểm tra thư viện phụ thuộc, kiểm tra hạ tầng, kiểm tra Docker image, kiểm tra ứng dụng sau khi triển khai và thu thập artifact làm bằng chứng.

Trong DocVault, DevSecOps được thể hiện qua Jenkins pipeline. Ngoài chạy test, pipeline còn chạy SAST, SCA, Trivy filesystem scan, Trivy image scan, Checkov IaC scan, Terraform validate, smoke test và DAST bằng OWASP ZAP. Những bước này bám theo tư tưởng ``shift-left security'' được đề cập trong các tài liệu DevSecOps \cite{owasp_devsecops,nist_ssdf}.

\section{Các nhóm kiểm thử bảo mật trong pipeline}

\begin{table}[H]
\centering
\caption{Các nhóm kiểm thử bảo mật trong DevSecOps pipeline}
\label{tab:security-test-types}
\begin{tabularx}{\textwidth}{p{0.14\textwidth}p{0.25\textwidth}Yp{0.22\textwidth}}
\toprule
\textbf{Nhóm} & \textbf{Mục tiêu} & \textbf{Ví dụ áp dụng trong DocVault} & \textbf{Công cụ} \\
\midrule
SAST & Phân tích mã nguồn tĩnh để phát hiện lỗi chất lượng và bảo mật. & Kiểm tra backend/frontend TypeScript, quality gate và code smell. & SonarQube \\
SCA & Phân tích thư viện phụ thuộc và lỗ hổng đã biết. & Quét dependency của monorepo Node.js. & OWASP Dependency Check \\
FS scan & Quét filesystem repository. & Phát hiện lỗ hổng trong file hệ thống, lockfile, cấu hình. & Trivy \\
Image scan & Quét Docker image sau build. & Phát hiện CVE trong image backend/frontend. & Trivy \\
IaC scan & Quét cấu hình hạ tầng. & Kiểm tra Terraform/Kubernetes manifest. & Checkov \\
DAST & Kiểm thử ứng dụng đang chạy. & Baseline scan web app sau deploy. & OWASP ZAP \\
\bottomrule
\end{tabularx}
\end{table}

\section{GitOps và Argo CD}

GitOps là mô hình trong đó Git đóng vai trò nguồn sự thật cho trạng thái triển khai. Khi pipeline build image mới, thay vì trực tiếp chỉnh cluster bằng lệnh thủ công, pipeline cập nhật Helm values hoặc manifest trong một branch GitOps. Argo CD theo dõi branch này và đồng bộ trạng thái mong muốn vào Kubernetes cluster \cite{argocd_docs}.

Trong DocVault, pipeline cập nhật các file values tương ứng với từng service. Sau đó Argo CD sync các application như gateway, metadata-service, document-service, workflow-service, audit-service, notification-service và web. Tách bạch như vậy giúp build pipeline và deployment controller giữ trách nhiệm riêng, đồng thời lưu lại lịch sử thay đổi triển khai trong Git.

\section{Tổng hợp công nghệ và công cụ}

Để có cái nhìn tổng thể, phần này tổng hợp các công nghệ và công cụ chính được sử dụng trong đề tài, chia thành hai nhóm: nhóm xây dựng ứng dụng (Bảng~\ref{tab:tech-stack-app}) và nhóm quy trình DevSecOps cùng hạ tầng triển khai (Bảng~\ref{tab:tech-stack-devsecops}). Các phiên bản được ghi nhận theo cấu hình thực tế trong repository tại thời điểm hoàn thành báo cáo.

\begin{table}[H]
\centering
\small
\caption{Công nghệ xây dựng ứng dụng DocVault}
\label{tab:tech-stack-app}
\begin{tabularx}{\textwidth}{p{0.20\textwidth}p{0.26\textwidth}p{0.14\textwidth}Y}
\toprule
\textbf{Thành phần} & \textbf{Công nghệ} & \textbf{Phiên bản} & \textbf{Vai trò} \\
\midrule
Frontend & Next.js, React & 16.2, 19.2 & Giao diện web theo App Router. \\
UI \& form & Tailwind CSS, Radix UI, React Hook Form, Zod & 4.x, 7.7, 4.3 & Thiết kế giao diện, form và validation phía client. \\
Server state & TanStack Query & 5.90 & Quản lý dữ liệu server, cache và đồng bộ. \\
Backend & NestJS & 10.0 & Framework cho gateway và các microservice. \\
Ngôn ngữ & TypeScript & 5.x & Ngôn ngữ chính cho frontend và backend. \\
Xác thực & Keycloak, passport-jwt, jwks-rsa & --- & OIDC/JWT, xác minh token bằng JWKS RS256. \\
ORM & Prisma + adapter-pg & 7.5 & Truy cập PostgreSQL bằng driver adapter. \\
Metadata DB & PostgreSQL & --- & Lưu metadata, ACL, workflow history. \\
Audit DB & MongoDB, Mongoose & --- & Lưu audit event và hash-chain. \\
Object storage & MinIO / Amazon S3 & --- & Lưu blob file tài liệu. \\
Rate limit & \code{@nestjs/throttler} & 6.5 & Giới hạn tần suất request theo cấu hình chung. \\
Monorepo & pnpm, Turbo & 9.15, 2.1 & Quản lý workspace và điều phối build. \\
\bottomrule
\end{tabularx}
\end{table}

\begin{table}[H]
\centering
\small
\caption{Công cụ DevSecOps và hạ tầng triển khai}
\label{tab:tech-stack-devsecops}
\begin{tabularx}{\textwidth}{p{0.22\textwidth}p{0.24\textwidth}Y}
\toprule
\textbf{Nhóm} & \textbf{Công cụ} & \textbf{Vai trò trong đề tài} \\
\midrule
CI/CD & Jenkins, Jenkins Shared Library & Điều phối pipeline, tách CI validation và CD release. \\
Secret scan & TruffleHog & Phát hiện credential bị commit trong repository. \\
SAST & SonarQube & Phân tích mã nguồn tĩnh và Quality Gate. \\
SCA & OWASP Dependency-Check & Quét lỗ hổng thư viện phụ thuộc theo CVSS. \\
FS/Image scan & Trivy & Quét filesystem và Docker image. \\
IaC scan & Checkov & Quét cấu hình Terraform, Kubernetes, Helm. \\
Policy-as-code & Kyverno CLI & Kiểm tra policy trên manifest đã render. \\
IaC & Terraform & Mô tả hạ tầng AWS dưới dạng mã. \\
Container & Docker BuildKit & Build và cache image theo batch. \\
Registry & Harbor & Private registry, lưu SBOM và trạng thái ký. \\
Ký artifact & Cosign & Ký và xác minh image digest. \\
GitOps & Git, Argo CD & Đồng bộ desired state theo mô hình app-of-apps. \\
Orchestration & Amazon EKS, Helm, Kustomize & Chạy workload trên Kubernetes. \\
Secret \& key & AWS Secrets Manager, External Secrets, KMS & Quản lý bí mật và khóa mã hóa (customer-managed key, xoay vòng hằng năm). \\
DAST & OWASP ZAP & Quét động ứng dụng sau triển khai. \\
Observability & Prometheus, Grafana, Loki & Thu thập metric, log và trực quan hóa. \\
\bottomrule
\end{tabularx}
\end{table}
